\documentclass[11pt,a4paper]{article}
\usepackage[margin=2.4cm]{geometry}
\usepackage[T1]{fontenc}
\usepackage[utf8]{inputenc}
\usepackage{lmodern}
\usepackage{textcomp}
\usepackage{booktabs}
\usepackage{longtable}
\usepackage{enumitem}
\usepackage{xcolor}
\usepackage{listings}
\usepackage{hyperref}
\hypersetup{colorlinks=true,linkcolor=blue!60!black,urlcolor=blue!60!black}
\setlist{nosep}
\lstset{basicstyle=\ttfamily\small,breaklines=true,columns=fullflexible,frame=single,
  framerule=0.3pt,rulecolor=\color{gray!50},backgroundcolor=\color{gray!6},
  xleftmargin=4pt,xrightmargin=4pt,aboveskip=6pt,belowskip=6pt}
\newcommand{\file}[1]{\texttt{#1}}
\newcommand{\param}[1]{\texttt{#1}}
\newcommand{\why}[1]{\par\smallskip\noindent\textit{Why this matters here:} #1\par}

\title{Gazebo and PX4 for the Atlas 10\,-\,EDF airframe\\
\large From nothing to a hover in SITL and HITL, with the tiltlab stack}
\author{tiltlab project notes}
\date{2026-09-11}

\begin{document}
\maketitle

\noindent This guide assumes you have never used Gazebo or PX4. It is written so that, with this
document and an internet connection, you can rebuild and operate everything the repository does
without further help. Where an official page explains something better than a rewrite would, the
guide sends you there and tells you what to take from it. Where this project learned something the
official pages do not say, the guide says it in full.

\tableofcontents
\newpage

%=====================================================================
\section{The map: what talks to what}

Everything here is three programs and one cable.

\begin{enumerate}
  \item \textbf{A physics simulator} (Gazebo) integrates rigid-body motion of the aircraft, spins
        simulated fans that produce thrust, and produces fake sensor readings.
  \item \textbf{A flight controller} (PX4) reads sensors, estimates attitude and position, runs the
        controllers and the \emph{control allocator} that turns ``roll torque, pitch torque, yaw
        torque, thrust'' into ten fan commands, and sends those commands back to the simulator.
  \item \textbf{tiltlab} (this repository) describes the aircraft (a \emph{scenario}: fan positions,
        fan thrust axes, mass, inertia, foil deflections), computes whether it can hover and how
        much control authority it has, and \emph{exports} the simulator model and the PX4
        configuration so that the two agree exactly.
  \item \textbf{The cable} is the difference between SITL and HITL:
    \begin{itemize}
      \item \textbf{SITL} (software in the loop): PX4 runs as a program on the same PC as Gazebo.
            No hardware. Fast to iterate.
      \item \textbf{HITL} (hardware in the loop): PX4 runs on the real Pixhawk, connected by USB.
            Gazebo replaces the aircraft and its sensors. The controller code, its timing and its
            parameters are exactly what will fly.
    \end{itemize}
\end{enumerate}

\begin{lstlisting}
scenario JSON -> tiltlab -> exported harness -> WSL launcher -> Gazebo <-> PX4 (PC or Pixhawk)
\end{lstlisting}

\why{Every failure in this project was a disagreement between two of those boxes: the model and
the allocator disagreed about geometry, the simulator and the board disagreed about attitude, or
the parameter file and the board disagreed about integers. Keep asking ``which two boxes disagree?''}

%=====================================================================
\section{Gazebo, the minimum you need}

\subsection{Two Gazebos}
There are two unrelated programs called Gazebo:
\begin{itemize}
  \item \textbf{Gazebo Classic} (versions 9, 10, 11; command \file{gazebo}). Old, stable, and the
        only one PX4 v1.17 speaks HITL through. We run it on Ubuntu 22.04.
  \item \textbf{gz sim} (formerly Ignition; versions Fortress, Garden, Harmonic; command \file{gz sim}).
        The current one. PX4's SITL bridge targets it. We run Harmonic on Ubuntu 24.04.
\end{itemize}
They have different file layouts, plugin systems and command-line tools. Do not mix their
documentation. Official entry points:
\begin{itemize}
  \item gz sim: \url{https://gazebosim.org/docs/harmonic/getstarted/} (read ``Getting started'' and
        ``Understanding the GUI''; skip ROS integration).
  \item Gazebo Classic: \url{https://classic.gazebosim.org/tutorials} (read ``Beginner: Overview'' and
        ``Build a robot: Model structure and requirements'').
\end{itemize}

\subsection{SDF: the file that describes everything}
Both Gazebos read \textbf{SDF} (Simulation Description Format), an XML dialect. Reference:
\url{http://sdformat.org/spec}. You need four concepts:
\begin{description}[style=nextline]
  \item[world] The scene: gravity, ground plane, light, which models to spawn and where. One file,
        \file{*.world} (Classic) or \file{*.sdf} (gz sim). Our exporter writes one per scenario.
  \item[model] The aircraft. A directory with \file{model.config} (name, author) and
        \file{model.sdf} (the content). Gazebo finds models through the \file{GAZEBO\_MODEL\_PATH}
        (Classic) or \file{GZ\_SIM\_RESOURCE\_PATH} (gz sim) environment variable; the launcher
        sets these.
  \item[link] A rigid body inside a model: mass, inertia, collision shape, visual shape, pose. Our
        model has \file{base\_link} (the airframe, mass and inertia from the scenario), ten
        \file{rotor\_N} links (one per fan) and an IMU link.
  \item[joint] Connects two links. Each rotor joint is a \emph{revolute} joint whose axis is the
        fan's spin axis.
\end{description}

Poses are ``x y z roll pitch yaw'' in metres and radians, in the \emph{parent's} frame. A rotor
link's pose therefore places the fan on the airframe and points it.

\subsection{Frames: the single most important paragraph}
Gazebo's body frame is \textbf{FLU}: x forward, y left, z up. PX4's is \textbf{FRD}: x forward, y
right, z down. Converting is a 180\textdegree{} rotation about x: $(x, y, z)_{\mathrm{FRD}} \to (x, -y,
-z)_{\mathrm{FLU}}$. Gazebo's world frame is \textbf{ENU} (x east, y north, z up); PX4's is
\textbf{NED}. A vehicle facing Gazebo's $+x$ (east) has PX4 heading 090\textdegree. That is not
a bug; you will see ``yaw 90'' on a level, still vehicle and it is correct.

tiltlab works in FRD internally and applies the FLU conversion once, when writing the model. If
you ever hand-edit a pose, do it in FLU.

\why{The airframe hovers at 25\textdegree{} nose-up. tiltlab's scenario field
\param{frame.hover\_pitch\_deg} rotates the airframe geometry into the flight-controller frame, so
the exported model's \file{base\_link} is level when the real airframe is at 25\textdegree. In
Gazebo you therefore \emph{see} the airframe mesh tilted 25\textdegree{} while PX4 reads ``level''.
That is the intended design, in both simulators.}

\subsection{Plugins}
A plugin is a shared library Gazebo loads for a model. Ours (Classic, from PX4's
\file{sitl\_gazebo-classic}):
\begin{itemize}
  \item \file{libgazebo\_motor\_model.so}, one per rotor. Thrust $= k_T\,\omega^2$ along the
        rotor link's $+z$, torque $= k_M \cdot$ thrust, with a spin-up time constant. The
        exporter sets $k_T$ from the fan curve and the time constant from the fan lag.
  \item \file{libgazebo\_imu\_plugin.so}, \file{\_magnetometer\_}, \file{\_barometer\_},
        \file{\_gps\_}: simulated sensors.
  \item \file{libgazebo\_mavlink\_interface.so}: the bridge. Reads sensors, sends them to PX4 as
        MAVLink; receives PX4's motor commands and publishes them to the motor plugins. Its
        parameters that matter: \param{serialEnabled}, \param{serialDevice} (HITL),
        \param{hil\_mode}, \param{hil\_state\_level}, \param{input\_scaling} per motor channel.
\end{itemize}
gz sim's equivalents live in PX4's \file{Tools/simulation/gz} and the \file{gz\_bridge} module
inside PX4 itself; PX4 subscribes to gz topics directly, no MAVLink in between.

\subsection{Topics and tools}
Gazebo publishes on internal topics. Useful commands (Classic):
\begin{lstlisting}
gz topic -l                      # list topics (note the world name in each path)
gz topic -e <topic>              # echo (fails silently on PX4's custom message types)
gz stats                         # real-time factor: keep it at 1.00 for HITL
gz world -w <world> --reset-all  # put every model back where it spawned
gz model -w <world> -m <name> -p # pose; -w is required unless the world is called "default"
\end{lstlisting}
gz sim: \file{gz topic -l}, \file{gz topic -e -t <topic>}, \file{gz model --list}.

%=====================================================================
\section{PX4, the minimum you need}

Official documentation: \url{https://docs.px4.io/v1.17/en/} (we pin v1.17.0). Read, in this
order: ``Basic Concepts'', ``Flight Modes (Multicopter)'', ``Parameters and Configurations'',
``Control Allocation'', ``Simulation'' overview, then ``Hardware in the Loop''.

\subsection{Parameters}
PX4 is configured by named parameters, e.g. \param{MC\_ROLLRATE\_P}. Each has a type: INT32 or
FLOAT. They live on the board (or in \file{parameters.bson} for SITL) and survive reboots once
saved. Three things this project learned the hard way:
\begin{enumerate}
  \item Over MAVLink, \textbf{INT32 parameters travel as raw bits in a float field}. Writing them
        as ordinary floats corrupts every integer parameter (\param{HIL\_ACT\_FUNC1} became
        1120534528, which is the bit pattern of $101.0f$), and reading them back the same wrong
        way agrees with itself, so nothing looks wrong. Write integers through the board's shell
        (\file{param set NAME VALUE}) or use \file{scripts/px4\_board.py}, which does that and
        decodes reads correctly.
  \item Some parameters are read \textbf{only at boot}: \param{SYS\_HITL}, \param{SYS\_AUTOSTART},
        \param{EKF2\_EN}, \param{SDLOG\_*}. Change, save, reboot, then check.
  \item An airframe file's \file{param set-default} lines never override a value the user saved;
        \file{param set} lines do. \file{1001\_rc\_quad\_x.hil} uses \file{set-default} for its
        4-rotor geometry, which is why our saved 10-rotor set survives it.
\end{enumerate}

\subsection{The control allocator}
PX4 does not know about ``quad'' or ``hexa''; it knows a matrix. For each rotor $i$ you give a
position \param{CA\_ROTORi\_PX/PY/PZ} (m, FRD, from the centre of gravity), a thrust axis
\param{CA\_ROTORi\_AX/AY/AZ} (unit vector, FRD), a thrust coefficient \param{CA\_ROTORi\_CT} and
a torque coefficient \param{CA\_ROTORi\_KM}. PX4 builds the effectiveness matrix (source:
\file{ActuatorEffectivenessRotors.cpp}), pseudo-inverts it and desaturates
(\file{ControlAllocationSequentialDesaturation.cpp}). tiltlab replicates that code line for line
so its predictions match what PX4 will do; the exported parameters are the same numbers the
replica used. Read ``Control Allocation'' in the docs once; then trust the exporter.

\subsection{Gains that matter for this airframe}
The rate loop (\param{MC\_ROLLRATE\_P/I/D} etc.), the attitude loop (\param{MC\_ROLL\_P}...),
the position loops (\param{MPC\_XY\_*}, \param{MPC\_Z\_*}), the hover throttle
\param{MPC\_THR\_HOVER}, the tilt limit \param{MPC\_TILTMAX\_AIR}, and \param{MC\_AT\_EN}
(autotune; \textbf{keep it 0}, autotune produced gains three to four times too high and every axis
rang). tiltlab sizes the first three from authority, inertia and fan lag; the rest were tuned in
SITL and saved in the scenario under \param{control.px4\_params\_override}, from where both
exports pick them up.

%=====================================================================
\section{Windows, WSL2 and USB}

Both simulators run in WSL2 (Windows Subsystem for Linux). Two distributions, because gz sim
Harmonic wants Ubuntu 24.04 and Gazebo Classic 11 exists only up to 22.04:
\begin{lstlisting}
wsl --list --verbose             # Ubuntu-24.04 (SITL), Ubuntu-22.04 (HITL)
wsl -d Ubuntu-22.04 -- bash -lc "<command>"
\end{lstlisting}
The repository is reachable inside both as \file{\textasciitilde/utopia/vibe-coded} (a symlink;
the space in ``Utopia Labs'' breaks quoting through \file{wsl.exe}, so never pass the
\file{/mnt/c/...} spelling on a command line).

\textbf{GUI}: WSLg shows Linux windows on the Windows desktop. If every window is grey and
titled \texttt{[WARN:COPY MODE]}, WSLg lost its shared-memory channel at start-up; only
\file{wsl --shutdown} from PowerShell fixes it (it closes every distribution).

\textbf{USB into WSL} uses usbipd-win (\url{https://github.com/dorssel/usbipd-win}):
\begin{lstlisting}
usbipd list                                   # find the Pixhawk's BUSID, e.g. 2-4
usbipd bind --busid 2-4                       # once, administrator PowerShell
usbipd attach --wsl Ubuntu-22.04 --busid 2-4  # each session (--auto-attach to survive reboots)
usbipd detach --busid 2-4                     # give it back to Windows (QGroundControl)
\end{lstlisting}
Inside WSL the board is \file{/dev/ttyACM0}, owned by \file{root:dialout}; your user must be in
\file{dialout}. A board reboot re-enumerates and often comes back as \file{ttyACM1}; the launcher
handles either. The distribution must be running for \file{attach} to work (keep a shell open).

%=====================================================================
\section{Installing the stack from scratch}

Do these once. Each step's check is written next to it.

\begin{enumerate}[leftmargin=*]
  \item \textbf{Windows tools}: \file{uv} (Python), Node, \file{make}, MiKTeX (for this guide),
        usbipd-win. In the repo: \file{source scripts/env.sh} (bash) or \file{. .\textbackslash
        scripts\textbackslash env.ps1} puts them on \file{PATH}. Check: \file{make test} is green.
  \item \textbf{WSL distributions}: \file{wsl --install -d Ubuntu-24.04} and
        \file{wsl --install -d Ubuntu-22.04}. Create the symlink in each:
        \file{mkdir -p \textasciitilde/utopia \&\& ln -s "/mnt/c/Users/<you>/Documents/Utopia Labs/vibe-coded" \textasciitilde/utopia/vibe-coded}.
  \item \textbf{PX4 source} in each distribution, pinned:
\begin{lstlisting}
git clone --recursive -b v1.17.0 https://github.com/PX4/PX4-Autopilot.git ~/PX4-Autopilot
cd ~/PX4-Autopilot && bash ./Tools/setup/ubuntu.sh     # toolchains; do NOT pass --no-nuttx in 22.04
\end{lstlisting}
        Check (22.04): \file{arm-none-eabi-gcc --version}. Check (24.04): \file{gz sim --version}
        prints Harmonic 8.x. In 22.04 also \file{sudo apt install gazebo libgazebo-dev}.
  \item \textbf{Let the launcher finish the rest}. Run
        \file{bash \textasciitilde/utopia/vibe-coded/scripts/wsl/tiltlab\_gazebo.sh --check}
        in each distribution. It lists every tool with its version, builds what is missing when
        asked, and applies the PX4 patches a 10-motor vehicle needs (PX4's gz bridge assumes at
        most 8 ESCs). Repeat until every line says \texttt{[ok]}.
  \item \textbf{HITL firmware} (22.04, once). Stock \file{px4\_fmu-v6x\_default} has no
        \file{pwm\_out\_sim} and adding it overflows flash by 22\,884 bytes, so the exporter ships
        a board label that also drops the fixed-wing and VTOL modules:
\begin{lstlisting}
bash ~/utopia/vibe-coded/scripts/wsl/tiltlab_gazebo.sh --build-firmware --harness <hitl harness dir>
\end{lstlisting}
        It copies \file{fmu-v6x\_hitl.px4board}, runs \file{make px4\_fmu-v6x\_hitl} (94.6\,\% of
        flash) and offers the upload. Check on the board: \file{pwm\_out\_sim status} lists
        channels 0..9 with functions 101..110 after the parameters are loaded.
\end{enumerate}

%=====================================================================
\section{Flying SITL}

\begin{enumerate}[leftmargin=*]
  \item In the app (\file{make dev}, \url{http://127.0.0.1:8000}) load or build a scenario, check
        \emph{hover: exact} and the authority numbers, Save.
  \item Metrics panel, \emph{Launch Gazebo}, \textbf{SITL}. The app exports
        \file{exports/gazebo/<name>/} (model, world, airframe \file{NNNN\_gz\_<name>}, README) and
        starts the launcher in Ubuntu-24.04. Its console tail appears under the buttons; the full
        console is \file{exports/logs/gazebo\_sitl.log}. Without the browser:
        \file{make menu}, option 2.
  \item Wait for \texttt{Ready for takeoff!} plus about ten seconds (height estimate), then
\begin{lstlisting}
wsl -d Ubuntu-24.04 -- bash ~/utopia/vibe-coded/scripts/wsl/px4ctl.sh takeoff
wsl -d Ubuntu-24.04 -- bash ~/utopia/vibe-coded/scripts/wsl/px4ctl.sh land
wsl -d Ubuntu-24.04 -- bash ~/utopia/vibe-coded/scripts/wsl/px4ctl.sh log   # copies the newest ulog
uv run --project backend python scripts/hover_report.py exports/logs/<file>.ulg
\end{lstlisting}
  \item Good hover on this airframe: attitude under 1\textdegree{} peak to peak, position under
        0.3\,m, torque demand well below 1. Tune live with \file{px4ctl.sh param NAME VALUE ...},
        land and take off again so integrators restart, and save winners in the scenario's
        \param{control.px4\_params\_override}.
\end{enumerate}

\textbf{What SITL cannot tell you}: PX4-SITL's clock \emph{is} Gazebo's clock (lockstep), so the
flight stack never sees late sensor data. The real board runs on its own clock. If a geometry
only just hovers in SITL, expect HITL to be worse.

%=====================================================================
\section{Flying HITL}

\subsection{How HITL differs, in this project's configuration}
The board is the flight controller; Gazebo Classic is the aircraft. Two designs exist for how the
board learns its attitude:
\begin{itemize}
  \item \param{hil\_state\_level 0}: Gazebo sends simulated IMU, magnetometer, barometer and GPS
        (\texttt{HIL\_SENSOR}, \texttt{HIL\_GPS}); the board's own EKF estimates attitude. The
        official default.
  \item \param{hil\_state\_level 1}: Gazebo sends its ground-truth attitude, position and velocity
        (\texttt{HIL\_STATE\_QUATERNION}); the board uses that directly.
\end{itemize}
This project uses level 1, and the exported parameter set carries what that needs:
\param{EKF2\_EN 0} (otherwise the EKF and the HIL handler both publish attitude and the values
flip between them), \param{SYS\_HAS\_MAG 0}, \param{SYS\_HAS\_BARO 0} (no such sensors are sent
at level 1), and calibration slot 0 pointed at the simulated IMU (device id 1310988) so the
arming checks stop looking for the real ICM. The reason for leaving level 0: the Classic IMU
plugin handed this board an accelerometer tilted about 25\textdegree{} while the model was
level, also with PX4's own \file{iris\_hitl} reference model, and the EKF believed it. That was
never resolved at the source; level 1 sidesteps it, and what HITL is for (the controller and
allocator running on real hardware against the real geometry) is fully exercised.

\subsection{Session checklist}
\begin{enumerate}[leftmargin=*]
  \item Fans and ESCs unpowered. Board on USB, attached to WSL (Section 4). QGroundControl closed
        (one program owns the serial port).
  \item Parameters onto the board, once per scenario (no sim running):
\begin{lstlisting}
python3 scripts/px4_board.py --dev /dev/ttyACM0 push   exports/gazebo_hitl/<name>_hitl/px4/<name>_hitl.params
python3 scripts/px4_board.py --dev /dev/ttyACM0 shell  reboot
python3 scripts/px4_board.py --dev /dev/ttyACM0 verify exports/gazebo_hitl/<name>_hitl/px4/<name>_hitl.params
\end{lstlisting}
        \file{verify} must print \texttt{N of N match}.
  \item App, \emph{Launch Gazebo}, \textbf{HITL} (or \file{make menu}, option 3). The launcher
        installs the model into the PX4 Classic tree, finds the serial device, starts the QGC
        relay (WSL2 drops UDP broadcast, so QGroundControl on Windows would otherwise never see
        the vehicle) and opens Gazebo.
  \item Before arming, confirm the board agrees with the sim:
\begin{lstlisting}
python3 scripts/px4_board.py status      # HIL flag True, roll/pitch near 0, velocities 0, Preflight check: OK
\end{lstlisting}
        (default link is the sim's SDK port 14540, so this works while Gazebo owns the serial.)
  \item Arm and fly from the transmitter or QGroundControl. Afterwards
        \file{python3 scripts/px4\_board.py pull-log exports/logs} then \file{hover\_report.py}.
        Logs land in \file{/fs/microsd/log/sessNNN/} because the board has no time source at level 1.
\end{enumerate}

\subsection{Returning the board to flight}
Reflash \file{px4\_fmu-v6x\_default}, load your flight parameter file, confirm \param{SYS\_HITL} 0,
\param{EKF2\_EN} 1, \param{SYS\_HAS\_MAG} 1, \param{SYS\_HAS\_BARO} 1, re-select your airframe
and recalibrate the IMU (the HITL set points calibration slot 0 at the simulated device).
On the real aircraft the IMU is bolted to the airframe and reads the 25\textdegree{} hover pitch;
\param{SENS\_BOARD\_Y\_OFF} (degrees) makes PX4's body frame the hover frame. It acts on real
sensors only, so it plays no part in HITL.

%=====================================================================
\section{When it goes wrong}

\begin{longtable}{p{4.6cm}p{4.9cm}p{5.4cm}}
\toprule
\textbf{You see} & \textbf{It means} & \textbf{Do} \\
\midrule
\endhead
Motors stay at zero after takeoff; allocator reports thrust unallocated & the geometry has no hover
trim (tiltlab shows \emph{Fz unattainable}) & pick a feasible sweep row or change hover pitch \\
Lifts, then flips within two seconds & PX4 default gains on a 12\,kg, 150\,ms-fan airframe & gains
from the exporter (\file{px4\_tuning}); never Autotune \\
Slow 0.25\,Hz roll/pitch swing to the tilt limit, metres of wander & position loop faster than
attitude loop & the exporter's \param{MPC\_*} scaling; lower \param{MPC\_TILTMAX\_AIR} on thin margin \\
Yaw torque demand above 1 & yaw axis has almost no authority; gains sized too stiff & soften
\param{MC\_YAWRATE\_P/I}, \param{MC\_YAW\_P} \\
Translates away at speed, altitude fine & a standing trim moment the integrator cannot hold &
raise \param{MC\_PR\_INT\_LIM} (pitch) with a larger \param{MC\_PITCHRATE\_I} \\
Second Gazebo launch shows a grey or duplicate window & a \file{gzserver} from the last run holds
port 11345 & \file{tiltlab\_gazebo.sh --stop} (the launcher also asks) \\
\texttt{[WARN:COPY MODE]} window & WSLg lost shared memory & \file{wsl --shutdown}, relaunch \\
\texttt{Error opening serial device: set\_option: Input/output error}, \texttt{Tx queue overflow}
& the board re-enumerated as another \file{ttyACM} & relaunch; the launcher picks the one present \\
Board attitude far from the sim's, parked; arms into a lunge & sim and board disagree about
attitude & \param{hil\_state\_level 1} + \param{EKF2\_EN 0} (in the export); compare with
\file{px4\_board.py status} before arming \\
\texttt{Preflight Fail: Accel Sensor 0 missing}, \texttt{barometer 0 missing}, \texttt{Found 0
compass} & sensor-presence checks at state level 1 & \param{SYS\_HAS\_MAG/BARO 0}, calibration
slot 0 = 1310988 (in the export) \\
\texttt{Flight termination active} & it flipped past 60\textdegree{}; latched & reboot the board
(\file{px4\_board.py shell reboot}); \param{CBRK\_FLIGHTTERM 121212} disables it for bench work \\
\texttt{logger: not running}, no log of the flight & \param{SDLOG\_BACKEND 0} on the board &
\param{SDLOG\_MODE 2}, \param{SDLOG\_BACKEND 1} (in the export), reboot \\
Integer parameters read as huge numbers or 0 after a MAVLink load & float write of INT32 &
\file{px4\_board.py push} (shell) and \file{verify} \\
QGroundControl ``Disconnected'' while Gazebo runs & WSL2 dropped the broadcast & the relay must be
running (launcher output); manual UDP link: listen 14551, server \texttt{<WSL ip>:18570} \\
\bottomrule
\end{longtable}

%=====================================================================
\section{Parameters used in this project}

Every parameter touched during the simulation and tuning phase, with what it does (from the
pinned PX4 v1.17.0 source: \file{PARAM\_DEFINE} doc blocks and \file{module.yaml}), the range or
enumeration PX4 accepts, and the value flown on the 50/55/60/65 geometry at 25\textdegree{} hover
pitch. Indexed families (\param{CA\_ROTORn\_*}, \param{HIL\_ACT\_FUNCn}, \param{SIM\_GZ\_EC\_*n})
are shown once with rotor 0's value. Values marked ``hand tuned'' live in the scenario's
\param{control.px4\_params\_override}; the rest are computed or fixed by the exporter.

{\small
\begin{longtable}{@{}p{3.3cm}p{5.6cm}p{3.6cm}p{1.9cm}@{}}
\toprule
\textbf{Parameter} & \textbf{What it does} & \textbf{Allowed} & \textbf{Used} \\
\midrule
\endfirsthead
\toprule
\textbf{Parameter} & \textbf{What it does} & \textbf{Allowed} & \textbf{Used} \\
\midrule
\endhead
\midrule
\multicolumn{4}{l}{\textbf{Geometry and allocation (from the scenario; identical in SITL and HITL)}} \\
\midrule
\param{CA\_AIRFRAME} & Airframe selection & 0 = Multirotor; 1 = Fixed-wing; 2 = Standard VTOL; 3 = Tiltrotor VTOL; 4 = Tailsitter VTOL; 5 = Rover (Ackermann); 6 = Rover (Differential); 7 = Motors (6DOF); 8 = Multirotor with Tilt; 9 = Custom; 10 = Helicopter (tail ESC); 11 = Helicopter (tail Servo); 12 = Helicopter (Coaxial); 13 = Rover (Mecanum); 14 = Spacecraft 2D; 15 = Spacecraft 3D & 0 \\
\param{CA\_METHOD} & Control allocation method & 0 = Pseudo-inverse with output clipping; 1 = Pseudo-inverse with sequential desaturation technique; 2 = Automatic & 2 \\
\param{CA\_ROTOR\_COUNT} & Total number of rotors & 0 = 0; 1 = 1; 2 = 2; 3 = 3; 4 = 4; 5 = 5; 6 = 6; 7 = 7; 8 = 8; 9 = 9; 10 = 10; 11 = 11; 12 = 12 & 10 \\
\param{CA\_ROTORn\_PX} & Position of rotor n along X body axis relative to center of gravity & -100 to 100 m & -0.0996811 \\
\param{CA\_ROTORn\_AX} & Axis of rotor n thrust vector, X body axis component & -100 to 100 & 0.258819 \\
\param{CA\_ROTORn\_CT} & Thrust coefficient of rotor n & 0 to 100 & 31.45 \\
\param{CA\_ROTORn\_KM} & Moment coefficient of rotor n & -1 to 1 & 0 \\
\param{THR\_MDL\_FAC} & Thrust to motor control signal model parameter & 0 to 1 & 1 \\
\midrule
\multicolumn{4}{l}{\textbf{Rate loop (sized by px4\_tuning; yaw and integrator limits hand tuned)}} \\
\midrule
\param{MC\_ROLLRATE\_P} & Roll rate P gain & 0.01 to 0.5 & 0.3898 \\
\param{MC\_ROLLRATE\_I} & Roll rate I gain & 0 to ... & 0.2339 \\
\param{MC\_ROLLRATE\_D} & Roll rate D gain & 0 to 0.01 & 0.01949 \\
\param{MC\_ROLLRATE\_K} & Roll rate controller gain & 0.01 to 5 & 1 \\
\param{MC\_RR\_INT\_LIM} & Roll rate integrator limit & 0 to ... & 0.3 \\
\param{MC\_PITCHRATE\_P} & Pitch rate P gain & 0.01 to 0.6 & 0.6 \\
\param{MC\_PITCHRATE\_I} & Pitch rate I gain & 0 to ... & 0.5 \\
\param{MC\_PITCHRATE\_D} & Pitch rate D gain & 0 to ... & 0.03 \\
\param{MC\_PR\_INT\_LIM} & Pitch rate integrator limit & 0 to ... & 0.6 \\
\param{MC\_YAWRATE\_P} & Yaw rate P gain & 0 to 0.6 & 0.08 \\
\param{MC\_YAWRATE\_I} & Yaw rate I gain & 0 to ... & 0.02 \\
\param{MC\_YAWRATE\_D} & Yaw rate D gain & 0 to ... & 0 \\
\param{MC\_YR\_INT\_LIM} & Yaw rate integrator limit & 0 to ... & 0.15 \\
\midrule
\multicolumn{4}{l}{\textbf{Attitude loop and autotune}} \\
\midrule
\param{MC\_ROLL\_P} & Roll P gain & 0 to 12 & 0.762 \\
\param{MC\_PITCH\_P} & Pitch P gain & 0 to 12 & 0.762 \\
\param{MC\_YAW\_P} & Yaw P gain & 0 to 5 & 0.25 \\
\param{MC\_AT\_EN} & Multicopter autotune module enable & 0 off, 1 on & 0 \\
\param{MC\_AT\_APPLY} & Controls when to apply the new gains & 0 = Do not apply the new gains (logging only); 1 = Apply the new gains after disarm; 2 = WARNING Apply the new gains in air & 0 \\
\midrule
\multicolumn{4}{l}{\textbf{Position and velocity loops, hover throttle, limits (hand tuned in SITL)}} \\
\midrule
\param{MPC\_THR\_HOVER} & Vertical thrust required to hover & 0.1 to 0.8 norm & 0.384 \\
\param{MPC\_XY\_P} & Proportional gain for horizontal position error & 0 to 2 & 0.25 \\
\param{MPC\_XY\_VEL\_P\_ACC} & Proportional gain for horizontal velocity error & 1.2 to 5 & 0.4 \\
\param{MPC\_XY\_VEL\_I\_ACC} & Integral gain for horizontal velocity error & 0 to 60 & 0.06 \\
\param{MPC\_XY\_VEL\_D\_ACC} & Differential gain for horizontal velocity error & 0.1 to 2 & 0.15 \\
\param{MPC\_XY\_VEL\_MAX} & Maximum horizontal velocity & 0 to 20 m/s & 1 \\
\param{MPC\_Z\_P} & Proportional gain for vertical position error & 0.1 to 1.5 & 0.4 \\
\param{MPC\_Z\_VEL\_P\_ACC} & Proportional gain for vertical velocity error & 2 to 15 & 1.2 \\
\param{MPC\_Z\_VEL\_I\_ACC} & Integral gain for vertical velocity error & 0.2 to 3 & 0.6 \\
\param{MPC\_Z\_VEL\_MAX\_UP} & Maximum ascent velocity & 0.5 to 8 m/s & 1 \\
\param{MPC\_Z\_VEL\_MAX\_DN} & Maximum descent velocity & 0.5 to 4 m/s & 0.7 \\
\param{MPC\_TILTMAX\_AIR} & Maximum tilt angle in air & 20 to 89 deg & 8 \\
\param{MPC\_TKO\_SPEED} & Takeoff climb rate & 1 to 5 m/s & 1 \\
\param{MPC\_ACC\_HOR} & Acceleration for autonomous and for manual modes & 2 to 15 m/s\^{}2 & 1 \\
\param{MPC\_ACC\_UP\_MAX} & Maximum upwards acceleration in climb rate controlled modes & 2 to 15 m/s\^{}2 & 2 \\
\param{MPC\_ACC\_DOWN\_MAX} & Maximum downwards acceleration in climb rate controlled modes & 2 to 15 m/s\^{}2 & 2 \\
\param{MPC\_JERK\_AUTO} & Jerk limit in autonomous modes & 1 to 80 m/s\^{}3 & 1 \\
\param{MIS\_TAKEOFF\_ALT} & Default take-off altitude & 0 to ... m & 2.5 (default) \\
\midrule
\multicolumn{4}{l}{\textbf{SITL only (gz sim bridge; written to the airframe file)}} \\
\midrule
\param{SIM\_GZ\_EC\_FUNCn} & Output function driven by gz ESC channel n (mixer\_module output functions) & 0 disabled; 101..112 Motor1..Motor12 & 101..110 \\
\param{SIM\_GZ\_EC\_MINn} & Command sent to gz at zero thrust (rotor rad/s) & 0 to 1000 & 0 \\
\param{SIM\_GZ\_EC\_MAXn} & Command sent to gz at full thrust (rotor rad/s) & 0 to 1000 & 1000 \\
\param{SIM\_GZ\_EN} & Simulator Gazebo bridge enable & 0 off, 1 on & 1 \\
\param{MAV\_0\_BROADCAST} & Broadcast heartbeats on local network for MAVLink instance n & 0 = Never broadcast; 1 = Always broadcast; 2 = Only multicast & 1 \\
\midrule
\multicolumn{4}{l}{\textbf{HITL only (written to the .params file loaded on the board)}} \\
\midrule
\param{SYS\_HITL} & Enable HITL/SIH mode on next boot & -1 = external HITL; 0 = HITL and SIH disabled; 1 = HITL enabled; 2 = SIH enabled & 1 \\
\param{SYS\_AUTOSTART} & Auto-start script index & airframe id; 1001 = HIL Quadcopter X & 1001 \\
\param{HIL\_ACT\_FUNCn} & Output function on pwm\_out\_sim channel n, sent in HIL\_ACTUATOR\_CONTROLS & 0 disabled; 101..112 Motor1..Motor12 & 101..110 \\
\param{CBRK\_SUPPLY\_CHK} & Circuit breaker for power supply check & 0 to 894281 & 894281 \\
\param{UAVCAN\_ENABLE} & UAVCAN mode & 0 = Disabled; 1 = Sensors Manual Config; 2 = Sensors Automatic Config; 3 = Sensors and Actuators (ESCs) Automatic Config & 0 \\
\param{EKF2\_EN} & EKF2 enable & 0 off, 1 on & 0 \\
\param{SYS\_HAS\_MAG} & Control if and how many magnetometers are expected & 0 off, 1 on & 0 \\
\param{SYS\_HAS\_BARO} & Control if the vehicle has a barometer & 0 off, 1 on & 0 \\
\param{CAL\_ACC0\_ID} & Accelerometer n calibration device ID & device id; 1310988 = the HIL receiver's SIM IMU & 1310988 \\
\param{CAL\_GYRO0\_ID} & Gyroscope n calibration device ID & device id; 1310988 = the HIL receiver's SIM IMU & 1310988 \\
\param{CAL\_ACC0\_*OFF} & Accelerometer offsets of slot 0 (X, Y, Z) & any (m/s\^{}2) & 0 \\
\param{CAL\_GYRO0\_*OFF} & Gyroscope offsets of slot 0 (X, Y, Z) & any (rad/s) & 0 \\
\param{CAL\_ACC0\_*SCALE} & Accelerometer scale of slot 0 (X, Y, Z) & 0.1 to 3 & 1 \\
\param{SDLOG\_MODE} & Logging Mode & 0 = when armed until disarm (default); 1 = from boot until disarm; 2 = from boot until shutdown; 3 = while manual input AUX1 \textgreater{}30\%; 4 = from 1st armed until shutdown & 2 \\
\param{SDLOG\_BACKEND} & Logging Backend (integer bitmask) & 0 to 3 & 1 \\
\param{GPS\_1\_CONFIG} & Serial port for GPS 1; rcS sets 0 in HITL (no GPS) & 0 disabled; a serial port id & 0 (set by rcS) \\
\midrule
\multicolumn{4}{l}{\textbf{Bench and real-aircraft settings mentioned in the text}} \\
\midrule
\param{CBRK\_FLIGHTTERM} & Circuit breaker for flight termination & 0 to 121212 & 0 (kept) \\
\param{SENS\_BOARD\_Y\_OFF} & Board rotation Y (pitch) offset & -45 to 45 deg & 0 in HITL; 25 on the real aircraft (sign to verify) \\
\param{SENS\_BOARD\_ROT} & Coarse board rotation relative to the body (multiples of 45/90 deg) & 0 = none; 1..40 = the rotation enum in PX4 docs (Sensor Orientation) & 0 \\
\bottomrule
\end{longtable}
}

%=====================================================================
\section{Where things live}

\begin{tabular}{@{}p{5.2cm}p{9.6cm}@{}}
\toprule
\file{scenarios/*.json} & aircraft descriptions; \param{control.px4\_params\_override} holds tuned gains \\
\file{backend/tiltlab/export/gazebo.py} & gz sim harness and gain sizing (\file{px4\_tuning}) \\
\file{backend/tiltlab/export/gazebo\_classic\_hitl.py} & Classic HITL harness, \file{.params}, board label, README \\
\file{backend/tiltlab/api/gazebo\_launch.py} & the app's launch/status/stop \\
\file{scripts/wsl/tiltlab\_gazebo.sh} & the launcher (checks, installs, patches, launches, \file{--stop}) \\
\file{scripts/wsl/px4ctl.sh} & SITL: takeoff, land, param, log, stop \\
\file{scripts/wsl/qgc\_udp\_relay.py} & MAVLink relay to QGroundControl on Windows \\
\file{scripts/px4\_board.py} & board: status, shell, push, verify, pull-log \\
\file{scripts/hover\_report.py} & hover quality from a ulog \\
\file{scripts/tiltlab\_menu.py} & \file{make menu}, the no-browser front door \\
\file{docs/sim\_stack.md} & the architecture map and the HITL findings \\
\file{docs/hitl\_runbook.md} & the HITL checklist for this machine \\
\file{exports/gazebo/}, \file{exports/gazebo\_hitl/} & generated harnesses (not versioned) \\
\file{third\_party/PX4-Autopilot} & pinned v1.17.0 source, cite file and line when porting \\
\bottomrule
\end{tabular}

\bigskip
\noindent\textbf{Further reading, in the order it pays off}: PX4 ``Basic Concepts'' and ``Control
Allocation''; gz sim ``Getting started''; sdformat.org spec (world, model, link, joint); PX4
``Simulation'' and ``Hardware in the Loop''; usbipd-win README; PX4 ``Parameters and
Configurations'' (the INT32 paragraph above is not in it).

\end{document}
